% Snapshot of the cross-references in P1_bernoulli_20260701_Bernoulli_revised.aux.
% Auto-regenerated on every compile; do not edit.
\companionlabel{sec:intro}{{1}{1}{Introduction}{section.1}{}}
\companionlabel{sec:theory}{{2}{4}{Methodology}{section.2}{}}
\companionlabel{sec:lemma}{{2.1}{4}{Equivalence of Forward and Tail Limits}{subsection.2.1}{}}
\companionlabel{lem:equiv}{{2.1}{4}{Equivalence of forward and tail limits}{theorem.2.1}{}}
\companionlabel{rem:interp-lemma}{{2.2}{5}{Interpretation for sequential learning}{theorem.2.2}{}}
\companionlabel{rem:nontail}{{2.3}{5}{Non-tail-measurable targets}{theorem.2.3}{}}
\companionlabel{ass:spectral}{{2.4}{5}{Effective operator-norm contractivity}{theorem.2.4}{}}
\companionlabel{sec:rnn-decomp}{{2.2}{5}{Martingale Decomposition of RNN Hidden States}{subsection.2.2}{}}
\companionlabel{thm:mds}{{2.5}{6}{MDS Decomposition with Stabilization Rate}{theorem.2.5}{}}
\companionlabel{eq:Mbound}{{1}{6}{MDS Decomposition with Stabilization Rate}{equation.2.1}{}}
\companionlabel{eq:Mstab}{{2}{6}{MDS Decomposition with Stabilization Rate}{equation.2.2}{}}
\companionlabel{eq:var-decomp}{{3}{6}{MDS Decomposition with Stabilization Rate}{equation.2.3}{}}
\companionlabel{eq:SigmaM}{{4}{6}{MDS Decomposition with Stabilization Rate}{equation.2.4}{}}
\companionlabel{eq:nondeg}{{5}{6}{MDS Decomposition with Stabilization Rate}{equation.2.5}{}}
\companionlabel{rem:rank}{{2.7}{7}{Scope of the nondegeneracy condition}{theorem.2.7}{}}
\companionlabel{cor:fclt}{{2.8}{7}{Functional CLT for MDS innovations}{theorem.2.8}{}}
\companionlabel{eq:fclt}{{6}{7}{Functional CLT for MDS innovations}{equation.2.6}{}}
\companionlabel{prop:predresid}{{2.9}{7}{Bayes prediction residuals are martingale differences}{theorem.2.9}{}}
\companionlabel{prop:rnn-resid-bridge}{{2.10}{8}{Bridge from recurrent states to prediction innovations}{theorem.2.10}{}}
\companionlabel{eq:bridge-bound}{{7}{8}{Bridge from recurrent states to prediction innovations}{equation.2.7}{}}
\companionlabel{eq:bridge-drift}{{8}{8}{Bridge from recurrent states to prediction innovations}{equation.2.8}{}}
\companionlabel{prop:whiteness}{{2.11}{8}{Innovation defect and the misspecification gap}{theorem.2.11}{}}
\companionlabel{rem:whiteness-arl}{{2.12}{8}{Consequence for calibration and the ARL guarantee}{theorem.2.12}{}}
\companionlabel{prop:whiteness-rate}{{2.13}{9}{Rate of the innovation defect and of the guarantee}{theorem.2.13}{}}
\companionlabel{sec:nonstationary}{{2.3}{9}{Extension to Non-Stationary Processes}{subsection.2.3}{}}
\companionlabel{def:ls}{{2.14}{9}{Locally stationary process; \citealt {dahlhaus1997fitting}}{theorem.2.14}{}}
\companionlabel{cor:ls}{{2.15}{9}{Approximate local stationary decomposition}{theorem.2.15}{}}
\companionlabel{eq:local-approx-bound}{{9}{10}{Approximate local stationary decomposition}{equation.2.9}{}}
\companionlabel{sec:pathway}{{3}{10}{Innovation CUSUM: In-Control False-Alarm Theory}{section.3}{}}
\companionlabel{eq:cusum}{{10}{10}{Setup and Assumptions}{equation.3.10}{}}
\companionlabel{ass:pathway}{{3.1}{10}{Innovation CUSUM operating conditions}{theorem.3.1}{}}
\companionlabel{thm:arl}{{3.2}{11}{Sharp in-control ARL lower bound for the Innovation CUSUM}{theorem.3.2}{}}
\companionlabel{eq:arl-lb}{{11}{11}{Sharp in-control ARL lower bound for the Innovation CUSUM}{equation.3.11}{}}
\companionlabel{rem:calibration}{{3.4}{11}{Calibration in practice}{theorem.3.4}{}}
\companionlabel{cor:twosided}{{3.5}{12}{Two-sided Innovation CUSUM}{theorem.3.5}{}}
\companionlabel{sec:delay}{{3.3}{12}{Out-of-Control Detection Delay}{subsection.3.3}{}}
\companionlabel{thm:delay}{{3.6}{12}{Out-of-control detection-delay upper bound}{theorem.3.6}{}}
\companionlabel{eq:delay}{{12}{12}{Out-of-control detection-delay upper bound}{equation.3.12}{}}
\companionlabel{cor:tradeoff}{{3.7}{12}{Exponential-ARL / logarithmic-delay trade-off}{theorem.3.7}{}}
\companionlabel{sec:optimality}{{3.4}{13}{Asymptotic Optimality}{subsection.3.4}{}}
\companionlabel{thm:optimality}{{3.9}{13}{Asymptotic optimality of the Innovation CUSUM}{theorem.3.9}{}}
\companionlabel{cor:dim}{{3.11}{14}{Dimension-stable ARL}{theorem.3.11}{}}
\companionlabel{eq:raw-bias}{{13}{14}{Dimension-stable ARL}{equation.3.13}{}}
\companionlabel{sec:empirical}{{4}{15}{Empirical Studies and Discussion}{section.4}{}}
\companionlabel{sec:processes}{{4.1}{15}{Processes, Benchmarks, and Real Data}{subsection.4.1}{}}
\companionlabel{sec:mds}{{4.2}{16}{Study MDS: Stationary Case --- Prediction-Residual MDS Diagnostic}{subsection.4.2}{}}
\companionlabel{tab:mds}{{1}{16}{Study MDS: Ljung--Box MDS pass rate (fraction of 1000 test sequences not rejecting zero autocorrelation at $\alpha = 0.05$, lag $\leq 5$). Bold values are within $\pm 0.05$ of the nominal level 0.95}{table.1}{}}
\companionlabel{sec:cusum-cp}{{4.3}{16}{Study CUSUM-CP: Innovation CUSUM at a Known Change-Point}{subsection.4.3}{}}
\companionlabel{tab:cusum-cp}{{2}{18}{Study CUSUM-CP: three-way CUSUM comparison at $h=5$. AR(1) change $\phi : 0.30 \to 0.85$ at $\tau ^* = 100$ ($T = 200$, $N = 2000$ sequences). Each replicate is classified by its first alarm into one of three mutually exclusive outcomes: a \emph {false positive} (FP, alarm before $\tau ^*$), a \emph {correct detection} (TP, first alarm at or after $\tau ^*$), or a \emph {miss} (FN, no alarm in the window); FP${}+{}$TP${}+{}$FN${}=1$. Mean delay$^\ddagger $ is the average of $(\text {alarm}-\tau ^*)$ over all replicates that alarm before the final step and therefore \emph {mixes} false positives (negative) with correct detections; it is a clean detection-speed measure only for methods whose FP rate is small. $\mathrm {ARL}_0$ is estimated on an independent in-control population and is the unconfounded false-positive metric. M3 (Innovation CUSUM) uses LSTM prediction residuals; no AR model is provided}{table.2}{}}
\companionlabel{fig:cusum-cp}{{1}{18}{Study CUSUM-CP: three-way CUSUM comparison. Left: false-positive rate (pre-$\tau ^*$ alarms; dashed line = nominal 5\%). Center: mean delay over any alarm (contaminated by false positives for Raw/Whitened; see Table~\ref {tab:cusum-cp}). Right: any-alarm rate (\emph {not} correct-detection rate --- for Raw this is dominated by pre-$\tau ^*$ false alarms). The Innovation CUSUM achieves the lowest false-positive rate and, per Table~\ref {tab:cusum-cp}, the highest correct-detection rate, without requiring knowledge of the AR model structure}{figure.1}{}}
\companionlabel{sec:cusum-dim}{{4.4}{18}{Study CUSUM-DIM: Innovation CUSUM Dimension Scaling}{subsection.4.4}{}}
\companionlabel{fig:cusum-dim-heat}{{2}{20}{Study CUSUM-DIM: ARL$_0$ (left) and false-alarm rate (right) as functions of dimension $d$ for the three CUSUM methods. The Innovation CUSUM (green) maintains a high, growing ARL$_0$ and a stable FA rate, while Raw (blue) and Whitened (orange) collapse as $d$ increases}{figure.2}{}}
\companionlabel{tab:cusum-dim}{{3}{20}{Study CUSUM-DIM: multivariate CUSUM comparison ($h=5$, 1000 test sequences, 2000 ARL$_0$ sequences). The two error types are reported separately: FP = false-positive rate (first alarm before $\tau $); TP = correct-detection rate (first alarm at or after $\tau $); FN = miss rate (no alarm in the window); FP${}+{}$TP${}+{}$FN${}=1$. Delay$^\ddagger $ = mean of $(\text {alarm}-\tau )$ over any alarm before the final step (negative = pre-$\tau $); it mixes false positives with correct detections and is a clean speed measure only when FP is small. Bold marks the best TP and ARL$_0$ per block}{table.3}{}}
\companionlabel{sec:etf}{{4.5}{20}{Study ETF: Real Multivariate Application --- US Equity Sector Returns}{subsection.4.5}{}}
\companionlabel{fig:cusum-dim-delay}{{3}{21}{Study CUSUM-DIM: Distribution of detection delay (steps from $\tau $ to alarm) across $d\in \{1,2,5,10\}$. The Innovation CUSUM (right column) consistently yields positive delays (valid post-change detection); Raw and Whitened (left and middle columns) shift to large negative delays as $d$ grows}{figure.3}{}}
\companionlabel{tab:etf}{{4}{22}{Study ETF: US equity sector ETF CUSUM results ($h=5$, $k=0.5$). Known change-point $\tau ^*=32$ (2020-02-19). ARL$_0$ estimated from 500 block-bootstrapped in-control score sequences of length~500, right-censored at the sequence length. Delay = alarm time $-$ $\tau ^*$ (negative = before the break)}{table.4}{}}
\companionlabel{sec:discussion}{{4.6}{22}{Summary and Interpretation}{subsection.4.6}{}}
\companionlabel{fig:etf}{{4}{23}{Study ETF: CUSUM trajectories on the COVID test window (2020-01-01 to 2022-06-30); the vertical dashed line marks $\tau ^*=32$ (2020-02-19). \emph {Left} (Part~A, $d=1$, XLF): the Innovation CUSUM (bottom) crosses threshold $h=5$ at $t=17$, fifteen steps before the official market peak. \emph {Right} (Part~B, $d=3$, XLF/XLK/XLE): the Innovation CUSUM again alarms at $t=17$, and its ARL$_0$ gap over M1/M2 widens relative to Part~A, consistent with the dimension-stability prediction of Corollary~\ref {cor:dim}}{figure.4}{}}
\companionlabel{sec:conclusion}{{5}{23}{Conclusion}{section.5}{}}
